% ===== UNIT 03: The ViT-B training recipe gap (1.2) =====
% Status: REWRITTEN  (change to REWRITTEN / REVIEWED when done; see CHECKLIST.md)
% Rules: style rewrite only. Keep every \cite, \ref, \label, number, equation and table value as is. Change the wording, not the content.
\subsection{The ViT-B training recipe gap}
The introduction of TransReID~\cite{he2021transreid} changed how people approached
vehicle ReID by applying a pure Vision Transformer
(ViT)~\cite{dosovitskiy2021vit} backbone to the task. TransReID augments the
standard ViT-B/16 architecture with two extra modules: the Side
Information Embedding (SIE), which encodes camera index and viewpoint as
learnable token offsets to reduce domain bias, and the Jigsaw Patch Module
(JPM), which shuffles patch tokens in a parallel branch to force the model to learn local
features~\cite{he2021transreid}. With a DeiT-B/16 backbone and SGD
(momentum 0.9, LR 0.008, cosine decay), TransReID reports 82.3\% mAP and 97.1\%
Rank-1 on VeRi-776~\cite{he2021transreid}.

Despite broad adoption of this architecture as a baseline, subsequent work has
primarily focused on augmenting it with new architectural
components: semantic alignment via language
models~\cite{li2023clipreid,lu2025clipsenet}, graph-based spatial
reasoning~\cite{shen2022hpgn}, or multi-branch CNN
structures~\cite{luo2023mbr4b}, instead of looking again at what a
properly tuned ViT-B training-and-post-processing pipeline leaves on the table.
No pure ViT-B method has been reported above 90\% mAP on VeRi-776. We push a
pure ViT-B/16 pipeline to 89.97\% mAP, approaching that mark, and, via
training-free fusion with a complementary CNN-semantic stream, reach 93.3\%.

